\documentclass[../../main.tex]{subfiles}
\graphicspath{{../../image/}{../../figure/}}
\begin{document}

\section{$BV$ 解}\label{section:ch10-bv解}

我们考虑严格双曲的守恒律方程组
\begin{align}
\partial_t U + \partial_x F(U) = 0.
\label{align:ch10-10-1-1}
\end{align}
在本章中，$U$ 恒表示 $(-\infty,\infty)\times(0,\infty)$ 上有界可测、属于 $BV_{\mathrm{loc}}$ 类的函数，并且是\eqref{align:ch10-10-1-1}的弱解。由\cref{chapter:ch01-平衡律}中 $BV$ 函数的理论，
\begin{align*}
(-\infty,\infty)\times(0,\infty) = \mathscr{C}\cup\mathscr{J}\cup\mathscr{I},
\end{align*}
其中 $\mathscr{C}$ 是 $U$ 的近似连续点集，$\mathscr{J}$ 表示 $U$ 的近似跳跃间断点（激波集）所成的集合，$\mathscr{I}$ 表示 $U$ 的不规则点集。$\mathscr{I}$ 的一维 Hausdorff 测度为零：$\mathscr{H}^{1}(\mathscr{I})=0$。激波集 $\mathscr{J}$ 本质上被（至多可数条）$C^{1}$ 弧的并集所覆盖。对任意 $(\bar{x},\bar{t})\in\mathscr{J}$，伴随着一侧近似极限 $U_{\pm}$ 以及一条斜率为 $s$（激波速度）的“切线”；正如\cref{chapter:ch01-平衡律}所示，它们由 Rankine--Hugoniot 跳跃条件相联系。


我们通篇假定\cref{chapter:ch08-容许激波}中引入的 Lax $E$ 条件在此以强形式成立：每一条激波都是压缩激波，但不是过压缩激波。也就是说，若 $U_{\pm}$ 为激波集上任一点对应的一侧极限，$s$ 为相应的激波速度，则存在 $i\in\{1,\cdots,n\}$ 使得
\begin{align}
\lambda_{i-1}(U_{\pm}) < \lambda_{i}(U_{+}) \leqslant s \leqslant \lambda_{i}(U_{-}) < \lambda_{i+1}(U_{\pm}).
\label{align:ch10-10-1-2}
\end{align}
在\eqref{align:ch10-10-1-2}中，当 $i=1$ 时第一个不等式不必出现，当 $i=n$ 时最后一个不等式不必出现。此外，由于\eqref{align:ch10-10-1-1}是严格双曲的，只要 $U$ 的振幅充分小，第一个与最后一个不等式便自动成立。


为方便起见，我们按\cref{chapter:ch01-平衡律}所述对 $U$ 作归一化。特别地，在每一点 $(\bar{x},\bar{t})\in\mathscr{C}$ 处，$U(\bar{x},\bar{t})$ 等于相应的近似极限 $U_{0}$。回顾 $\mathscr{H}^{1}(\mathscr{J})=0$ 并利用\cref{chapter:ch01-平衡律}，容易推知存在零测集 $\mathcal{N}\subset(0,\infty)$，具有如下性质：对任一固定 $\bar{t}\notin\mathcal{N}$，函数 $U(\cdot,\bar{t})$ 在 $(-\infty,\infty)$ 上局部有界变差；$(\bar{x},\bar{t})\in\mathscr{C}$ 当且仅当 $U(\bar{x}-,\bar{t})=U(\bar{x}+,\bar{t})$，而 $(\bar{x},\bar{t})\in\mathscr{J}$ 当且仅当 $U(\bar{x}-,\bar{t})\neq U(\bar{x}+,\bar{t})$。在后一情形，$U_{-}=U(\bar{x}-,\bar{t})$，$U_{+}=U(\bar{x}+,\bar{t})$。


上述 $U$ 的性质仅仅来源于 $U\in BV$。而 $U$ 同时是\eqref{align:ch10-10-1-1}的解这一事实，应当诱导出额外的结构。基于后续章节将要讨论的特殊方程组的经验，可以合理地预期以下事实：$U$ 应在 $\mathscr{C}$ 上（古典地）连续，且 $\mathscr{J}$ 上点处的一侧极限 $U_{\pm}$ 应在古典意义下取到。此外，$\mathscr{I}$ 应是构成 $\mathscr{J}$ 的那些弧的（至多可数个）端点所成的集合。关于 $\mathscr{I}$ 中任一点作 $(x,t)$ 坐标的一致伸缩，在极限情形应得到\cref{chapter:ch09-容许波扇与riemann问题}中所述性质的波扇，亦即 $\mathscr{I}$ 应由激波生成点与激波相互作用点构成。上述图景在多大程度上准确刻画一般双曲守恒律方程组解的结构，将在后续章节中讨论。


\end{document}